% GENERATED FILE. Edit canonical lessons, then run npm run generate:books.
% canonical-source-hash: fnv1a64:c933bfc2c873116b
% canonical-lessons: TA-C60-needle, TA-C60-rope, TA-C60-umbrella, TA-C60-broom, TA-C60-comb

\chapter{Things People Make}
\label{ch:ta-things-people-make}
\hlchaptermodality{60}

\begin{chapteropening}
\textbf{By the end of this chapter:} \emph{I can name a needle, a rope, an umbrella, a broom and a comb in Tamil.}

Complete TA-C60-comb and say all five words in order.
\end{chapteropening}

\section[\texorpdfstring{\ta{ஊசி} (ūsi)}{ஊசி (ūsi)}]{\ta{ஊசி} (ūsi) — a needle}

\label{lesson:TA-C60-needle}



\begin{quote}
Before the new one: what did \emph{vellam} mean, and what did \emph{sarkkarai} mean?
\end{quote}

\subsection*{You'll want to know: \ta{ஊசி}}
\textbf{\ta{ஊசி}} (\emph{ūsi}) — \textquotedblleft{}a needle\textquotedblright{}.

The sewing needle, and whatever else is that shape. It is the tool \ta{துணி} is repaired with, which is most of what \ta{துணி} asks for.

Kannada \ka{ಸೂಜಿ} (\emph{sūji}) and Telugu \te{సూది} (\emph{sūdi}) hold on to a sound at the front that Tamil let go of, so Tamil's is the shortest of the three. Open it on a long \ta{ஊ}.

\textbf{In the family, and next door}

\noindent
\begin{tabularx}{\linewidth}{@{}>{\raggedright\arraybackslash}X>{\raggedright\arraybackslash}X>{\raggedright\arraybackslash}X@{}}
\toprule
\textbf{} & \textbf{word} & \textbf{same root} \\
\midrule
Kannada & \ka{ಸೂಜಿ} (\emph{sūji}) &  \\
Telugu & \te{సూది} (\emph{sūdi}) &  \\
Malayalam & \ml{സൂചി} (\emph{sūci}) &  \\
Hindi (neighbour) & \dv{सूई} (\emph{sūī}) & yes \\
English & a needle &  \\
\bottomrule
\end{tabularx}

\begin{grammarlens}[title={What you've built}]
The first of five made things.
\end{grammarlens}

\subsection*{Your turn}
\begin{itemize}
\raggedright
  \item \emph{Say it:} \emph{ūsi}
  \item \emph{Say it:} it once more, slowly
  \item \emph{Say it:} \emph{ūsi}, then \emph{tuṇi}, and say which one repairs the other
  \item \emph{Recall:} say \emph{pasu}
  \item \emph{Read it:} \textbf{\ta{தயிர்}}
\end{itemize}

\begin{tcolorbox}[breakable,colback=teal!4,colframe=teal!35!black,title={Before you move on}]
What does \ta{ஊசி} mean? (\textquotedblleft{}a needle\textquotedblright{}.) Say it once more.
\end{tcolorbox}
\section[\texorpdfstring{\ta{கயிறு} (kayiṟu)}{கயிறு (kayiṟu)}]{\ta{கயிறு} (kayiṟu) — a rope}

\label{lesson:TA-C60-rope}



\begin{quote}
Before the new one: what did \emph{ūsi} mean?
\end{quote}

\subsection*{You'll want to know: \ta{கயிறு}}
\textbf{\ta{கயிறு}} (\emph{kayiṟu}) — \textquotedblleft{}a rope\textquotedblright{}.

Twisted up out of coconut fibre, which is what a coast with that many palms was going to make rope out of. Thin ones tie a bundle; thick ones draw water.

The \ta{ற} at the end is the hard one — the same \ta{ற} standing in \ta{ஆறு} and in \ta{பறவை}. A \ta{பசு} in a yard is at the end of a \ta{கயிறு}, which is the commonest place to meet the word.

\textbf{In the family, and next door}

\noindent
\begin{tabularx}{\linewidth}{@{}>{\raggedright\arraybackslash}X>{\raggedright\arraybackslash}X>{\raggedright\arraybackslash}X@{}}
\toprule
\textbf{} & \textbf{word} & \textbf{same root} \\
\midrule
Kannada & \ka{ಹಗ್ಗ} (\emph{hagga}) &  \\
Telugu & \te{తాడు} (\emph{tāḍu}) &  \\
Malayalam & \ml{കയർ} (\emph{kayaṟ}) & yes \\
Hindi (neighbour) & \dv{रस्सी} (\emph{rassī}) &  \\
English & a rope &  \\
\bottomrule
\end{tabularx}

\begin{grammarlens}[title={What you've built}]
Two.
\end{grammarlens}

\subsection*{Your turn}
\begin{itemize}
\raggedright
  \item \emph{Say it:} \emph{kayiṟu}
  \item \emph{Say it:} it once more, slowly
  \item \emph{Say it:} \emph{kayiṟu}, then \emph{pasu}, and say which one holds the other
  \item \emph{Read it:} \textbf{\ta{ஆடு}}
  \item \emph{Recall:} say \emph{ney}
\end{itemize}

\begin{tcolorbox}[breakable,colback=teal!4,colframe=teal!35!black,title={Before you move on}]
What does \ta{கயிறு} mean? (\textquotedblleft{}a rope\textquotedblright{}.) Say it once more.
\end{tcolorbox}
\section[\texorpdfstring{\ta{குடை} (kuṭai)}{குடை (kuṭai)}]{\ta{குடை} (kuṭai) — an umbrella}

\label{lesson:TA-C60-umbrella}



\begin{quote}
Before the new one: what did \emph{kayiṟu} mean?
\end{quote}

\subsection*{You'll want to know: \ta{குடை}}
\textbf{\ta{குடை}} (\emph{kuṭai}) — \textquotedblleft{}an umbrella\textquotedblright{}.

Against \ta{மழை} and against the \ta{சூரியன்} equally, and in Tamil country the second job is the larger one. It is carried in a month with no rain in it without anybody thinking that odd.

Kannada \ka{ಕೊಡೆ} (\emph{koḍe}) is the same word. It ends on the \ta{ை} you already put on \ta{கூடை} and on \ta{விதை}.

\textbf{In the family, and next door}

\noindent
\begin{tabularx}{\linewidth}{@{}>{\raggedright\arraybackslash}X>{\raggedright\arraybackslash}X>{\raggedright\arraybackslash}X@{}}
\toprule
\textbf{} & \textbf{word} & \textbf{same root} \\
\midrule
Kannada & \ka{ಕೊಡೆ} (\emph{koḍe}) &  \\
Telugu & \te{గొడుగు} (\emph{goḍugu}) &  \\
Malayalam & \ml{കുട} (\emph{kuṭa}) & yes \\
Hindi (neighbour) & \dv{छाता} (\emph{chātā}) &  \\
English & an umbrella &  \\
\bottomrule
\end{tabularx}

\begin{grammarlens}[title={What you've built}]
Three.
\end{grammarlens}

\subsection*{Your turn}
\begin{itemize}
\raggedright
  \item \emph{Say it:} \emph{kuṭai}
  \item \emph{Say it:} it once more, slowly
  \item \emph{Say it:} \emph{kuṭai}, then \emph{maḻai}, and say which one you take out for the other
  \item \emph{Recall:} say \emph{kōḻi}
  \item \emph{Read it:} \textbf{\ta{எண்ணெய்}}
\end{itemize}

\begin{tcolorbox}[breakable,colback=teal!4,colframe=teal!35!black,title={Before you move on}]
What does \ta{குடை} mean? (\textquotedblleft{}an umbrella\textquotedblright{}.) Say it once more.
\end{tcolorbox}
\section[\texorpdfstring{\ta{துடைப்பம்} (tuṭaippam)}{துடைப்பம் (tuṭaippam)}]{\ta{துடைப்பம்} (tuṭaippam) — a broom}

\label{lesson:TA-C60-broom}



\begin{quote}
Before the new one: what did \emph{kuṭai} mean?
\end{quote}

\subsection*{You'll want to know: \ta{துடைப்பம்}}
\textbf{\ta{துடைப்பம்}} (\emph{tuṭaippam}) — \textquotedblleft{}a broom\textquotedblright{}.

Coconut ribs bound at one end, held short, swung low. It is the first thing that touches a \ta{வாசல்} in the morning: the ground is swept, then wetted, and only then does the \ta{கோலம்} go down on it.

Three houses, three words again — Kannada \ka{ಪೊರಕೆ} (\emph{porake}), Telugu \te{చీపురు} (\emph{cīpuru}), Tamil \ta{துடைப்பம்}. The doubled \ta{ப்ப} in the middle is the one you have been holding since \ta{பெட்டி} taught you to hold a letter twice.

\textbf{In the family, and next door}

\noindent
\begin{tabularx}{\linewidth}{@{}>{\raggedright\arraybackslash}X>{\raggedright\arraybackslash}X@{}}
\toprule
\textbf{} & \textbf{word} \\
\midrule
Kannada & \ka{ಪೊರಕೆ} (\emph{porake}) \\
Telugu & \te{చీపురు} (\emph{cīpuru}) \\
Malayalam & \ml{ചൂല്} (\emph{cūlŭ}) \\
Hindi (neighbour) & \dv{झाड़ू} (\emph{jhāṛū}) \\
English & a broom \\
\bottomrule
\end{tabularx}

\begin{grammarlens}[title={What you've built}]
Four.
\end{grammarlens}

\subsection*{Your turn}
\begin{itemize}
\raggedright
  \item \emph{Say it:} \emph{tuṭaippam}
  \item \emph{Say it:} it once more, slowly
  \item \emph{Say it:} \emph{tuṭaippam}, then \emph{kōlam}, and say which one comes first
  \item \emph{Read it:} \textbf{\ta{காக்கை}}
  \item \emph{Recall:} say \emph{vellam}
\end{itemize}

\begin{tcolorbox}[breakable,colback=teal!4,colframe=teal!35!black,title={Before you move on}]
What does \ta{துடைப்பம்} mean? (\textquotedblleft{}a broom\textquotedblright{}.) Say it once more.
\end{tcolorbox}
\section[\texorpdfstring{\ta{சீப்பு} (sīppu)}{சீப்பு (sīppu)}]{\ta{சீப்பு} (sīppu) — a comb}

\label{lesson:TA-C60-comb}



\begin{quote}
Before the new one: what did \emph{tuṭaippam} mean?
\end{quote}

\subsection*{You'll want to know: \ta{சீப்பு}}
\textbf{\ta{சீப்பு}} (\emph{sīppu}) — \textquotedblleft{}a comb\textquotedblright{}.

A comb, and nothing borrowed about it. The \ta{ீ} at the front is held long, and then the doubled \ta{ப்ப} comes again, so it lands in the same shape as \ta{துடைப்பம்} and \ta{பெட்டி}.

A \ta{சீப்பு} and a \ta{மலர்} are what go into hair on a morning that is being made something of, and the two turn up in the same sentence often enough to be worth learning together.

\textbf{In the family, and next door}

\noindent
\begin{tabularx}{\linewidth}{@{}>{\raggedright\arraybackslash}X>{\raggedright\arraybackslash}X>{\raggedright\arraybackslash}X@{}}
\toprule
\textbf{} & \textbf{word} & \textbf{same root} \\
\midrule
Kannada & \ka{ಬಾಚಣಿಗೆ} (\emph{bācaṇige}) &  \\
Telugu & \te{దువ్వెన} (\emph{duvvena}) &  \\
Malayalam & \ml{ചീപ്പ്} (\emph{cīppŭ}) & yes \\
Hindi (neighbour) & \dv{कंघी} (\emph{kaṅghī}) &  \\
English & a comb &  \\
\bottomrule
\end{tabularx}

\begin{grammarlens}[title={What you've built}]
Five, and the run is closed: a needle, a rope, an umbrella, a broom, and a comb.
\end{grammarlens}

\subsection*{Your turn}
\begin{itemize}
\raggedright
  \item \emph{Say it:} \emph{sīppu}
  \item \emph{Say it:} it once more, slowly
  \item \emph{Say it:} all five in order — \emph{ūsi}, \emph{kayiṟu}, \emph{kuṭai}, \emph{tuṭaippam}, \emph{sīppu}
  \item \emph{Recall:} say \emph{paṟavai}
  \item \emph{Read it:} \textbf{\ta{சர்க்கரை}}
\end{itemize}

\begin{tcolorbox}[breakable,colback=teal!4,colframe=teal!35!black,title={Before you move on}]
What does \ta{சீப்பு} mean? (\textquotedblleft{}a comb\textquotedblright{}.) Say it once more.
\end{tcolorbox}
